% ============================================================================
% COMMUNICATION PROTOCOL ANALYSIS
% ============================================================================
\section{Communication Protocol Analysis}

\subsection{The Problem Statement}

The STAR robot requires communication across three distinct paths:

\begin{enumerate}
    \item \textbf{Controller App $\rightarrow$ RPi5 Gateway}: Over WiFi/Internet (this is the focus)
    \item \textbf{RPi5 Gateway $\rightarrow$ ESP32}: Over SPI bus (already defined)
    \item \textbf{RPi5 Gateway $\rightarrow$ LiDAR}: Over UART (already defined)
\end{enumerate}

The critical question is: \textit{What protocol should Path \#1 use?}

\subsection{Protocol Options Evaluated}

Four options were evaluated for controller-to-robot communication:

\begin{enumerate}
    \item gRPC + GraphQL Hybrid
    \item ConnectRPC (modern gRPC alternative)
    \item WebSocket + GraphQL
    \item Pure gRPC with Envoy Proxy
\end{enumerate}

\subsection{Detailed Comparison}

\begin{table}[H]
\centering
\small
\begin{tabularx}{\textwidth}{|l|X|X|X|X|}
\hline
\textbf{Aspect} & \textbf{gRPC + GraphQL} & \textbf{ConnectRPC} & \textbf{WebSocket + GraphQL} & \textbf{Pure gRPC} \\
\hline
Latency & gRPC: 1-5ms, GraphQL: 10-30ms & 5-15ms & WS: 5-10ms, GQL: 10-30ms & 1-5ms \\
\hline
Packet Size & Binary (gRPC) + JSON (GQL) & Binary OR JSON & JSON only & Binary only \\
\hline
Browser Support & gRPC-Web (limited) + GQL (full) & Full (HTTP/1.1) & Full native & Requires Envoy \\
\hline
Bidirectional & gRPC-Web: NO, GQL Sub: YES & NO & YES (native) & gRPC-Web: NO \\
\hline
Infrastructure & 2 servers & 1 server & 1 server + WS & 1 server + Envoy \\
\hline
Learning Curve & High (2 systems) & Medium & Medium & Medium-High \\
\hline
Ecosystem & Very mature (both) & Newer (2022+) & Very mature & Very mature \\
\hline
\end{tabularx}
\caption{Protocol Comparison Matrix}
\end{table}

\subsection{Protocol Deep Dive}

\subsubsection{gRPC + GraphQL Hybrid (Selected)}

\begin{figure}[H]
\centering
\begin{tikzpicture}[node distance=2cm, auto]
    \node[draw, rectangle, rounded corners, fill=green!20, minimum width=3cm] (browser) {Controller App (Browser)};

    \node[draw, rectangle, rounded corners, fill=blue!20, minimum width=3cm, below left=1.5cm and 1cm of browser] (graphql) {GraphQL (port 8080)};
    \node[draw, rectangle, rounded corners, fill=red!20, minimum width=3cm, below right=1.5cm and 1cm of browser] (grpc) {gRPC-Web (port 8090)};

    \node[draw, rectangle, rounded corners, fill=orange!20, minimum width=4cm, below=3cm of browser] (gateway) {RPi5 Gateway};

    \draw[->, thick] (browser) -- (graphql) node[midway, left, font=\small] {Telemetry, Settings};
    \draw[->, thick] (browser) -- (grpc) node[midway, right, font=\small] {Motor Commands};
    \draw[->, thick] (graphql) -- (gateway);
    \draw[->, thick] (grpc) -- (gateway);
\end{tikzpicture}
\caption{gRPC + GraphQL Hybrid Architecture}
\end{figure}

\textbf{Advantages:}
\begin{itemize}
    \item gRPC binary format provides fastest motor command transmission
    \item GraphQL enables flexible queries for complex UI data needs
    \item Clear separation of concerns: real-time control vs. data queries
    \item Both protocols are industry standards with excellent tooling
\end{itemize}

\textbf{Disadvantages:}
\begin{itemize}
    \item Two different APIs to maintain
    \item gRPC-Web requires Envoy proxy as translator
    \item More complex client code (two client libraries)
\end{itemize}

\subsubsection{Why This Was Selected for STAR}

\begin{enumerate}
    \item \textbf{Low-Latency Motor Control}: gRPC's binary Protobuf format minimizes serialization overhead for time-critical motor commands.

    \item \textbf{Flexible Telemetry Queries}: GraphQL subscriptions provide real-time telemetry streaming with selective field fetching.

    \item \textbf{Offline Capability}: Apollo Client (GraphQL) has excellent caching and offline support built-in.

    \item \textbf{Latency Acceptability}: For motor commands, even 10-30ms latency from browser to RPi5 is acceptable because:
    \begin{itemize}
        \item The ESP32 maintains autonomous 1kHz PID control between commands
        \item Human reaction time is approximately 200-300ms
        \item The critical real-time path is RPi5 $\rightarrow$ ESP32 (SPI), not Controller $\rightarrow$ RPi5
    \end{itemize}
\end{enumerate}

\subsection{Packet Format Comparison}

\subsubsection{Protobuf (gRPC) Binary Encoding}

Protobuf uses a compact binary format that excludes field names:

\begin{lstlisting}[language=Protobuf, caption=Protobuf Message Definition]
message VelocityCommand {
    float left_velocity = 1;   // 4 bytes + tag
    float right_velocity = 2;  // 4 bytes + tag
    uint64 timestamp = 3;      // 1-10 bytes (varint)
}
\end{lstlisting}

\textbf{Wire format example} (VelocityCommand with left=1.0, right=0.5, timestamp=1702500000000):
\begin{verbatim}
0D 00 00 80 3F  // field 1 (float): 1.0
15 00 00 00 3F  // field 2 (float): 0.5
18 80 A0 CF EE E5 31  // field 3 (varint): timestamp
\end{verbatim}

\textbf{Total size: approximately 17 bytes}

\subsubsection{JSON (GraphQL) Text Encoding}

Equivalent JSON payload:
\begin{lstlisting}[language=json, caption=JSON Equivalent]
{
  "leftVelocity": 1.0,
  "rightVelocity": 0.5,
  "timestamp": 1702500000000
}
\end{lstlisting}

\textbf{Total size: approximately 65 bytes} (3.8x larger)

\subsubsection{Serialization Performance}

Based on industry benchmarks\footnote{Source: Ambassador Labs, Protobuf vs JSON Performance Analysis, 2024}:

\begin{table}[H]
\centering
\begin{tabular}{|l|c|c|c|}
\hline
\textbf{Metric} & \textbf{Protobuf} & \textbf{JSON} & \textbf{Difference} \\
\hline
Message Size & 17 bytes & 65 bytes & 3.8x smaller \\
Serialization (JVM) & 1.2 $\mu$s & 4.5 $\mu$s & 3.75x faster \\
Deserialization (JVM) & 0.8 $\mu$s & 3.2 $\mu$s & 4x faster \\
Network Latency (1Mbps) & 0.14 ms & 0.52 ms & 3.7x faster \\
\hline
\end{tabular}
\caption{Protobuf vs JSON Performance Comparison}
\end{table}
